\documentclass[12pt]{article}
\usepackage[utf8]{inputenc}
\usepackage[T1]{fontenc}
\usepackage{hyperref}
\title{Katkari Language Resources for K-12 Curriculum Development}
\author{Compiled by Claude Sonnet 5.5 (claude-sonnet-5-5) for Andr\'as Kornai}
\date{}
\begin{document}
\maketitle

\section{Introduction}
Katkari (Kathodi; ISO 639-3 kfu) is an Indo-Aryan language of Maharashtra, reported with about 12,000 native speakers (2007) and described as endangered. Katkari people commonly use Marathi as well. Only academic and wiki-type material was found.

\section{Primary Educational Resources}
None found.

\section{Linguistic and Reference Materials}
\begin{itemize}
\item S.\ B.\ Kulkarni, \emph{Descriptive analysis of K\=atkar\=i dialect}, Deccan College, Poona, 1969 (typed dissertation, 696 pp.), catalogued by Glottolog: \url{https://glottolog.org/resource/reference/id/hh:gtd:Kulkarni:Katkari}
\item Community-built Katkari--English dictionary on Glosbe (crowdsourced, unverified): \url{https://glosbe.com/kfu/en}
\item Language article (stub): \url{https://en.wikipedia.org/wiki/Katkari_language}
\item Article on the Katkari people: \url{https://en.wikipedia.org/wiki/Katkari_people}
\end{itemize}

\section{Religious and Community Resources}
None found.

\section{Digital and Audio Resources}
Only the Glosbe dictionary above.

\section{Curriculum Development Context}
No school materials were found. The Kulkarni dissertation is the most substantial source identified; access likely requires a library.
\end{document}